% TOP_PROBLEM: {"id": 1250, "title": "Quasiperfect Numbers", "category_id": 1, "category": {"id": 1, "name": "number_theory", "display_name": "Number Theory", "description": "Properties of integers, prime numbers, Diophantine equations.", "slug": "number-theory", "order_index": 1, "created_at": "2026-07-31T15:26:25.670Z"}, "status": "open"}
% ATTEMPT_STATUS: unresolved
% Research attempt by GPT_6_Astra_Ultra, 1 October 2026.
% Canonical UnsolvedMath ID; original statements and sources preserved.
% FIRST_POSTED: 2026-10-01
% Imported from: unsolvedmath_additional_500.tex; UnsolvedMath ID 1250
% !TeX program = lualatex
% Compile twice: lualatex 1250.tex
% Self-contained: no external bibliography, images, or \input files.
% Numeric section labels and visible numbers are original UnsolvedMath IDs.
\documentclass[11pt,a4paper]{article}
\usepackage[margin=24mm,headheight=14pt]{geometry}
\usepackage{fontspec}
\setmainfont{Latin Modern Roman}
\setsansfont{Latin Modern Sans}
\setmonofont{Latin Modern Mono}
\usepackage{amsmath,amssymb,mathtools}
\usepackage{unicode-math}
\setmathfont{Latin Modern Math}
\usepackage{microtype}
\usepackage{enumitem,xurl,needspace}
\usepackage[dvipsnames]{xcolor}
\usepackage{hyperref}
\hypersetup{unicode=true,colorlinks=true,linkcolor=MidnightBlue,urlcolor=MidnightBlue,
 pdftitle={UnsolvedMath Problem 1250},pdfauthor={GPT 6 Astra Ultra},bookmarksnumbered=true}
\usepackage{bookmark,tocloft,titlesec,fancyhdr}
\setlength{\cftsecnumwidth}{4.2em}
\cftsetpnumwidth{2.5em}
\cftsetrmarg{4em}
\titleformat{\section}{\Large\bfseries\raggedright}{\thesection}{0.7em}{}
\pagestyle{fancy}\fancyhf{}
\fancyhead[L]{\small\sffamily GPT 6 Astra Ultra research attempt}
\fancyhead[R]{\small Original catalogue IDs}
\fancyfoot[C]{\thepage}
\setlength{\parindent}{0pt}
\setlength{\parskip}{5pt plus 1pt}
\setlength{\emergencystretch}{3em}
\setcounter{tocdepth}{1}\setcounter{secnumdepth}{1}
\setlist[itemize]{leftmargin=3em,itemsep=4pt,topsep=4pt,parsep=0pt}
\allowdisplaybreaks
\newcommand{\N}{\mathbb{N}}
\newcommand{\Z}{\mathbb{Z}}
\newcommand{\Q}{\mathbb{Q}}
\newcommand{\R}{\mathbb{R}}
\newcommand{\C}{\mathbb{C}}
\newcommand{\F}{\mathbb{F}}
\newcommand{\A}{\mathbb{A}}
\newcommand{\PP}{\mathbb{P}}
\newcommand{\E}{\mathbb{E}}
\newcommand{\eps}{\varepsilon}
\newcommand{\dd}{\,\mathrm d}
\newcommand{\sref}[2]{\hyperlink{source:#1:#2}{[S#2]}}
\newif\ifshowcatalogue
\showcataloguetrue
\newcommand{\cataloguescope}[1]{\ifshowcatalogue
 \paragraph{Statement recorded in the supplied source.}
 #1\fi}
\begin{document}
% UnsolvedMath ID 1250; source catalogue code NT-043

\setcounter{section}{1249}

\section{Existence of Quasiperfect Numbers}\label{1250}

{\small\sffamily UnsolvedMath ID 1250 · Number Theory}

\subsection{Definitions and mathematical statement}

\paragraph{Shared notation.}Unless a section says otherwise, $\N=\{1,2,\ldots\}$,
$\Z$ is the ring of integers, $\Q,\R,\C$ are the rational, real, and complex
fields, $[N]=\{1,\ldots,\lfloor N\rfloor\}$, and $\log$ is the natural
logarithm. For real functions of a parameter tending to infinity,
$f\ll g$ and $f=O(g)$ mean $|f|\leq Cg$ eventually for some fixed $C$;
$f\gg g$ reverses this comparison; $f=o(g)$ means $f/g\to0$;
$f\sim g$ means $f/g\to1$; and $f\asymp g$ means both order bounds.
A subscript on $O$ or $\ll$ specifies allowed constant dependence.
The expression $n^{o(1)}$ in an upper bound means growth smaller than
$n^\epsilon$ for each fixed $\epsilon>0$. Local definitions govern any
exception, finite-field convention, or use of the same symbol for another
object. An asymptotic claim is not automatically an effective algorithm.



Write $\sigma(n)=\sum_{d\mid n}d$, summing over positive divisors. A quasiperfect number is a positive integer with $\sigma(n)=2n+1$, equivalently with the sum of its proper positive divisors equal to $n+1$. The problem asks whether any such integer exists. \sref{1250}{1} \sref{1250}{2}

\paragraph{Statement recorded in the supplied source.}
Do quasiperfect numbers exist? \sref{1250}{1}

\paragraph{Residual task reported by the archive.}

The embedded review (2026-08-17) records: Construct a quasiperfect number or prove none exist; independently audit and peer-review the new eight-prime-factor computation. \sref{1250}{1}

\subsection{Short English statement}

Can the proper divisors of a positive integer add up to exactly one more than the integer itself?

\subsection{Research attempt}

\paragraph{Scope and source status.}
The equation is $\sigma(n)=2n+1$, with the plus sign essential. The
2026 Toyohara--Tao--Yao preprint [S2] claims an eight-distinct-prime
lower bound using extensive computation. Its arXiv abstract was
inspected on 1 October 2026, but neither its full proof nor its code and
formal artifacts were audited here. That claimed bound is not used as
a premise. The attempt instead derives an elementary odd-square
reduction and a precise obstruction for a prescribed prime support.

\paragraph{Odd divisor sum forces an odd square part.}
For an odd prime $p$, $\sigma(p^e)$ is odd exactly when $e$ is even,
because it is a sum of $e+1$ odd terms. The factor
$\sigma(2^a)=2^{a+1}-1$ is always odd. Since $2n+1$ is odd, a
candidate must therefore have the shape
\[
                         n=2^a m^2,
 \qquad a\ge0,\quad m\text{ odd}.
\]
This parity argument alone does not imply that $n$ is odd. The even
case requires a separate obstruction, provided next.

\paragraph{The even case is impossible.}
Assume $a\ge1$ and set $A=2^{a+1}$, so $A\equiv0\pmod4$.
Multiplicativity gives
\[
 (A-1)\sigma(m^2)=Am^2+1,
 \qquad
 (A-1)\bigl(\sigma(m^2)-m^2\bigr)=m^2+1.
\]
Thus $A-1$ divides $m^2+1$. But $A-1\equiv3\pmod4$, so its
prime factorization contains a prime $r\equiv3\pmod4$ to an odd
exponent. In particular $r\mid m^2+1$, which is impossible: it would
give $m^2\equiv-1\pmod r$, while raising to the odd power
$(r-1)/2$ contradicts Fermat's theorem. Here $r\nmid m$ follows
immediately from $m^2\equiv-1$.
Therefore $a=0$ and every quasiperfect candidate is an odd square.
The case $m=1$ also fails directly, since $\sigma(1)=1\ne3$.
This proof tracks the even case instead of assuming a parity conclusion
stronger than the parity calculation provides.

\paragraph{Coprimality and support constraints.}
The defining equation implies
\[
             \gcd(n,\sigma(n))=\gcd(n,2n+1)=1.
\]
Consequently for any exact prime power $p^{2e}\Vert n$, every prime
divisor of $\sigma(p^{2e})$ is absent from $n$. This is the opposite
support behavior to odd perfect numbers, where divisor sums feed back
into the prime support. Mixing those two propagation rules would make
a factor-chain argument incorrect.

For instance, if $3^2\Vert n$ then $13\nmid n$, because
$\sigma(3^2)=13$. If $13^2\Vert n$, then neither $3$ nor $61$
can divide $n$, since $\sigma(13^2)=3\cdot61$. Thus the exact
exponent-two branch cannot contain both three and thirteen.
The word ``exact'' matters: changing the exponent changes the
divisor-sum polynomial and its forbidden primes.

Abundancy gives a complementary lower bound on the support size.
Since $n$ is odd,
\[
 2<\frac{\sigma(n)}n=\prod_{p^{2e}\Vert n}
       \left(1+p^{-1}+\cdots+p^{-2e}\right)
   <\prod_{p\mid n}\frac p{p-1}.
\]
With at most two distinct odd primes the last product is at most
$(3/2)(5/4)=15/8$, impossible. Hence at least three distinct primes
are necessary by this elementary argument. This is deliberately weaker
than the literature's much deeper reported bounds.

\paragraph{A strict reciprocal identity and the last-exponent problem.}
Dividing $\sigma(n)=2n+1$ by $n$ and pairing divisors gives
\[
                 \sum_{d\mid n}\frac1d=2+\frac1n.
\]
The excess over two is exactly $1/n$, not an arbitrary small positive
number. A proposed prime support with limiting abundancy at most two is
therefore impossible. If its limiting product exceeds two, however,
increasing the exponents may move the finite product across two, and
the exact excess $1/n$ still needs to be matched. Approximate closeness
of abundancy to two does not settle that equality.

There is an exact way to expose the remaining exponent on a fixed
branch. Write $n=Mp^{2e}$ with $(M,p)=1$, and put $S=\sigma(M)$.
Using the geometric series and multiplying by $p-1$ gives
\[
 S(p^{2e+1}-1)=(p-1)(2Mp^{2e}+1),
\]
or
\[
 \bigl(pS-2M(p-1)\bigr)p^{2e}=S+p-1.
\]
For fixed $M,p$, the coefficient on the left must be positive and
the positive rational quotient
\[
             \frac{S+p-1}{pS-2M(p-1)}
\]
must be an even power of $p$ with exponent at least two.
If the coefficient is zero or negative, the branch is impossible.
This is a complete test for one prescribed cofactor and last prime,
not a finite test for all possible $M$ and $p$.

\paragraph{Executed finite check.}
The script calculates $\sigma(n)$ exactly for every positive
$n\le200000$ using its prime factorization. None satisfies
$\sigma(n)=2n+1$. This independent check is small and is not a
reproduction of the 2026 preprint's claimed computation. It also
provides the exact small divisor-sum factorizations used in the support
examples. The files are \texttt{checks.py} and \texttt{results.json}
under \path{_Local/Erdos_problems/UnsolvedMath/attempt_audit_2026_10_01/number_theory/}.
No claim of proof-assistant verification or literature-code auditing
is made by this local arithmetic check.

\paragraph{Remaining gap and outcome.}
The attempt excludes every even integer and all odd supports of at
most two primes, derives the forbidden-support rule, and gives an exact
last-exponent equation on fixed branches. It does not limit the number
of branches or exclude odd squares with larger support.
\textbf{Outcome: UNRESOLVED.} The catalogue's request to independently
audit the recent large computation is not completed by the present
small search and is explicitly left separate from the proved reductions.

\subsection{Sources}

{\small

\begin{itemize}

\item[\hypertarget{source:1250:1}{S1}] ULAM.AI, \emph{UnsolvedMath}, supplied \texttt{problems.json}, record 1250 (NT-043), statement and background. The embedded literature-review date is 2026-08-17; that is the archive’s review date, not a fresh certification. Dataset landing page: \url{https://huggingface.co/datasets/ulamai/UnsolvedMath}.

\item[\hypertarget{source:1250:2}{S2}] Akira Toyohara, Ye Tao, and Siqiong Yao, Every quasiperfect number has at least eight distinct prime factors, arXiv:2608.02066 (2026). \url{https://arxiv.org/abs/2608.02066}. Bibliographic information imported from the supplied record.

\item[\hypertarget{source:1250:3}{S3}] Peter Hagis Jr. and Graeme L. Cohen, Some results concerning quasiperfect numbers, Journal of the Australian Mathematical Society, Series A 33 (1982), 275-286. \url{https://doi.org/10.1017/S1446788700018790}. Bibliographic information imported from the supplied record.


\item[E1] Akira Toyohara, Ye Tao, and Siqiong Yao,
\emph{Every quasiperfect number has at least eight distinct prime
factors}, arXiv:2608.02066v1 (2026).
\url{https://arxiv.org/abs/2608.02066v1}. Abstract inspected
1 October 2026. Its computational and formal-verification claims are
reported as claims; they were not independently reproduced here.

\end{itemize}

}

\label{end:1250}
\end{document}
